% Includable fragment. Requires amsmath, amssymb, amsthm and the
% theorem/lemma/proposition/corollary environments. Labels use rev:.
% Written proof; the accompanying suggested Lean module has NOT been compiled.

\section{Reversal makes the source restriction elementary}
\label{rev:section}

This section replaces the current kernel's equational transport by a
uniform crossed elimination. Its purpose is to simplify the equality
argument while preserving exactly the native consequence relation. The
proof uses no arithmetic induction: a single, syntactically nonzero
retraction context supplies everything needed to remove the restriction
on the source of an equality.

\subsection{Complete calculus and ordinary reading}

Numerical terms and native formulas are unchanged:
\[
 t::=x\mid0\mid S(t)\mid I(n,a;u.h),
 \qquad
 A::=(s=t)\mid Q_yA\mid D_{x,y}A.
\]
Only $u$ is bound in the body $h$ of an iterator. The binders $Q$ and $D$
bind respectively $y$ and $(x,y)$ in their matrices. All substitutions
below are capture-avoiding; unused template variables are permitted.
The ordinary readings are
\[
 Q_yB(y)\equiv B(0)\to\forall y\,B(y),\qquad
 D_{x,y}B(x,y)\equiv
 \exists x\bigl(B(x,0)\to\forall y\,B(x,y)\bigr).
\]
The designated contradiction remains $0=S0$.

Define an operation on the root of a formula by
\[
 (s=t)^\dagger=(t=s),\qquad
 (Q_yB)^\dagger=Q_yB,\qquad
 (D_{x,y}B)^\dagger=D_{x,y}B.
 \tag{Reversal}\label{rev:rho}
\]
It does not recurse into a compound matrix. In particular,
$(A^\dagger)^\dagger=A$, and reversal commutes with capture-avoiding
numerical substitution. Neither statement expresses an object-theory
axiom: these are identities of finite syntax.

Write $s\not\equiv0$ to mean that $s$ is not the literal zero term.
This condition does not assert $s\ne0$ in arithmetic. In particular,
$I(0,0;u.0)\not\equiv0$, although its value is zero.

The calculus $\mathcal R$ has exactly the following eight rule schemas.
The retained context $\Gamma$ is a finite reusable list of native
formulas.
\[
 \frac{}{\Gamma\vdash A}\ (A\in\Gamma),\qquad
 \frac{}{\Gamma\vdash I(0,a;u.h)=a},\qquad
 \frac{}{\Gamma\vdash I(Sn,a;u.h)=h[I(n,a;u.h)/u]}.
 \tag{Hyp, $I_0$, $I_S$}\label{rev:initial}
\]
\[
 \frac{\Gamma\vdash s=t\qquad\Gamma\vdash F(s)}
      {\Gamma\vdash F(t)^\dagger}
 \quad(s\not\equiv0),
 \tag{X}\label{rev:x}
\]
where $F$ is any native formula template, including a constant template,
and
\[
 \frac{\Gamma\vdash Q_yB(y)\qquad\Gamma\vdash B(0)}
      {\Gamma\vdash B(t)^\dagger}.
 \tag{QX}\label{rev:qx}
\]
The remaining three rules are exactly those of the current native kernel:
\[
 \frac{\Gamma\vdash D_{x,y}B(x,y)
       \qquad\Gamma,Q_yB(b,y)\vdash A}
      {\Gamma\vdash A},
 \tag{DE}\label{rev:de}
\]
\[
 \frac{\Gamma,B(t,0)\vdash B(t,Sz)}
      {\Gamma\vdash D_{x,y}B(x,y)},
 \tag{DI-case}\label{rev:di}
\]
\[
 \frac{\Gamma,E(z)\vdash E(Sz)}
      {\Gamma\vdash Q_yE(y)}
 \quad(E\text{ a literal equation}).
 \tag{QI-progress}\label{rev:qi}
\]
In DE, $b$ is fresh for $\Gamma$, $A$, and the fixed parameters of $B$.
In DI-case and QI-progress, $z$ is fresh for the retained context and
fixed parameters; in DI-case it is also fresh for the chosen witness
$t$. The displayed local assumption is discharged. No logical axiom,
reflexivity axiom, cut rule, unrestricted substitution rule, numerical
case rule, separation axiom, or additional induction rule is primitive.
Their admissibility, when needed, follows from the equivalence proved
below. In particular, no rule-by-rule substitution into a derivation of
$\mathcal R$ is assumed in this proof.

\subsection{Reversal and the equality algebra}

Throughout this subsection, all displayed derivations are over an
arbitrary fixed context $\Gamma$. Introduce the abbreviations
\[
             J(a):=I(0,a;u.0),\qquad c:=J(0).
 \label{rev:j}
\]
The equation $J(a)=a$ is an instance of $I_0$. Both $J(a)$ and $c$
are syntactically nonzero, independently of the value of $a$.

\begin{lemma}[Unrestricted reversal]\label{rev:reverse}
From $\Gamma\vdash A$ one can derive $\Gamma\vdash A^\dagger$,
for every native formula $A$.
\end{lemma}
\begin{proof}
Apply X with major premise $c=0$, supplied by $I_0$, and the constant
template $F(v):=A$. Its minor premise is the given $A$, and its
conclusion is $A^\dagger$. The source condition holds because
$c\not\equiv0$. This use of a constant template is essential.
\end{proof}

\begin{corollary}[Uncrossed rules]\label{rev:uncrossed}
Ordinary QE is admissible. Ordinary transport
$s=t,F(s)\Rightarrow F(t)$ is admissible whenever $s\not\equiv0$.
\end{corollary}
\begin{proof}
Apply QX or X, respectively, and then Lemma~\ref{rev:reverse}.
The two reversals cancel syntactically.
\end{proof}

\begin{lemma}[Reflexivity and symmetry]\label{rev:refl-symm}
Every $a=a$ is derivable. From $a=b$ one can derive $b=a$, with no
restriction on either endpoint.
\end{lemma}
\begin{proof}
For reflexivity use $J(a)=a$ twice, as the major and minor premises of X,
with template $F(v):=(v=a)$. The conclusion is $(a=a)^\dagger=a=a$.
Symmetry is the equational case of Lemma~\ref{rev:reverse}. Notice that
this includes equations with the literal source $0$.
\end{proof}

\begin{lemma}[Joining at zero]\label{rev:zero-join}
If $p,q\not\equiv0$, then from $p=0$ and $q=0$ one can derive $p=q$.
\end{lemma}
\begin{proof}
Use major premise $p=0$, minor premise $J(p)=p$, and template
$F(v):=(J(v)=p)$ in X. Its conclusion is
\[
                     (J(0)=p)^\dagger=(p=c).
\]
Similarly obtain $q=c$. Reverse both equations to get $c=p$ and
$c=q$. Finally apply X with major $c=q$, minor $c=p$, and template
$F(v):=(v=p)$. Its conclusion is $(q=p)^\dagger=(p=q)$.
Every crossed source used here is one of $p,q,c$, all syntactically
nonzero. No substitution into a proof, induction, or numerical case
distinction occurs.
\end{proof}

\begin{lemma}[Unrestricted transitivity]\label{rev:trans}
From $r=s$ and $s=t$ one can derive $r=t$.
\end{lemma}
\begin{proof}
If $s\not\equiv0$, apply the uncrossed restricted transport of
Corollary~\ref{rev:uncrossed} to major $s=t$, minor $r=s$, and
template $F(v):=(r=v)$.

Suppose instead that $s\equiv0$. If $r\equiv0$, the second premise is
already the conclusion. If $t\equiv0$, the first premise is already
the conclusion. In the remaining case $r,t\not\equiv0$, the premises
are $r=0$ and $0=t$. Reverse the latter and apply
Lemma~\ref{rev:zero-join}.

The distinction between a literal zero and another term is a decidable
case distinction on finite syntax in the metaproof. It is not an
application of excluded middle or induction in the object calculus.
\end{proof}

\begin{lemma}[Unrestricted numerical congruence]\label{rev:congr}
For every numerical template $f(v)$, from $s=t$ one can derive
$f(s)=f(t)$.
\end{lemma}
\begin{proof}
If $s\not\equiv0$, apply uncrossed restricted transport to major $s=t$
and template $F(v):=(f(s)=f(v))$, starting from reflexivity
$f(s)=f(s)$.

If $s\equiv0$ and $t\not\equiv0$, reverse the major to $t=0$ and use
template $F(v):=(f(v)=f(t))$. Reflexivity $f(t)=f(t)$ is the required
minor premise, so the conclusion is $f(0)=f(t)$.

If both endpoints are literally zero, the required equation is an
instance of reflexivity.
\end{proof}

\begin{theorem}[Recovery of equational transport]\label{rev:tr-eq}
The rule
\[
 \frac{\Gamma\vdash s=t\qquad\Gamma\vdash a(s)=b(s)}
      {\Gamma\vdash a(t)=b(t)}
\]
is admissible in $\mathcal R$, without a restriction on $s$.
\end{theorem}
\begin{proof}
Numerical congruence gives $a(s)=a(t)$ and $b(s)=b(t)$. Reverse the
first equation and use unrestricted transitivity twice:
\[
                         a(t)=a(s)=b(s)=b(t).
\]
\end{proof}

\subsection{The abstract retraction argument}

The preceding proof has a stronger, arithmetic-independent form. This
form identifies precisely which facts are used, and is also a concise
way to check that the equality argument is not circular.

\begin{proposition}[Retraction criterion]\label{rev:abstract}
Suppose a term language has a distinguished literal $0$ and a fixed
one-variable term context $J(v)$ satisfying both of the following
conditions:
\begin{enumerate}
\item $J(a)$ is syntactically different from $0$ for every term $a$;
\item $J(a)=a$ is an available axiom schema, uniformly in the retained
context.
\end{enumerate}
For equations, let reversal exchange their endpoints. Then X with
sources restricted to terms other than literal $0$ derives all ordinary
equational transport. The derivation uses only the retraction schema
and X.
\end{proposition}
\begin{proof}
Set $c=J(0)$. A single use of X with constant template $A$ and major
$c=0$ reverses any equation $A$. The retraction instance $J(a)=a$ used
twice derives reflexivity. For a nonzero term $p$ satisfying $p=0$,
the template $J(v)=p$ turns its retraction instance into $p=c$.
Thus any two nonzero terms equal to $0$ can be joined through the
nonzero anchor $c$. Ordinary transitivity now follows by splitting on
whether its middle term is literally $0$. Congruence follows by using
the nonzero endpoint of the source equation, reversing its orientation
when necessary; the equation $0=0$ is handled by reflexivity. Combining
congruence, symmetry, and transitivity gives general equational
transport, as above.
\end{proof}

For $\mathcal R$, take $J(v)=I(0,v;u.0)$. Consequently the entire
equality normalization theorem uses neither $I_S$, nor QX, nor DE,
nor either introduction rule. In particular, the numerical induction
rule is not invoked even indirectly. This contrasts with the supplied
ordinary source-restricted transport proof, whose zero-source symmetry
construction uses numerical tests and equation induction. The present
argument applies to a changed crossed rule, rather than giving a
shorter proof of that older rule's difficult intermediate lemma.

\subsection{Exact equivalence with the current kernel}

Let $\mathcal N$ denote the native kernel formalized as
\texttt{NPrf} at the pinned source commit. Its equality rule is
equational transport, its closure elimination is ordinary QE, and its
remaining six rules are Hyp, $I_0$, $I_S$, DE, DI-case, and QI-progress.

\begin{theorem}[Identity-language consequence equivalence]
\label{rev:equivalence}
For every finite native context $\Gamma$ and formula $A$,
\[
                 \Gamma\vdash_{\mathcal R}A
                 \quad\Longleftrightarrow\quad
                 \Gamma\vdash_{\mathcal N}A.
\]
There are effective, syntax-preserving translations of finite
derivations in both directions.
\end{theorem}
\begin{proof}
For $\mathcal R\to\mathcal N$, translate each finite derivation by
induction on its construction. The six unchanged rules translate
directly. For X, use the existing derived theorem
\texttt{NPrf.tr}, which transports through an arbitrary native
template, and then reverse the resulting equality if the formula is
atomic. If its root is $Q$ or $D$, reversal is the identity. For QX,
use ordinary QE followed by the same reversal. Equality symmetry is an
existing derived rule of $\mathcal N$.

For $\mathcal N\to\mathcal R$, again recurse on the derivation. The
six common rules are unchanged, equational transport is supplied by
Theorem~\ref{rev:tr-eq}, and ordinary QE is supplied by
Corollary~\ref{rev:uncrossed}. These constructions work over arbitrary
retained contexts, so they apply unchanged inside the fresh-variable
subderivations of DE and the two introduction rules.

There is no need to reprove compound transport inside $\mathcal R$:
the old kernel requires only its equational case in this direction.
The existing compound-transport theorem is used only for the
$\mathcal R\to\mathcal N$ translation, where it is already available.
\end{proof}

\begin{corollary}[Bi-interpretability and equiconsistency]
\label{rev:biinterpretation}
The existing native interpretations into and out of full first-order
iPRA are also a bi-interpretation between $\mathcal R$ and iPRA.
The numerical domains remain whole, equality retains its ordinary
meaning, and both interpretation composites have the same identity
graph $x=y$ as before. In particular, $\mathcal R$ is equiconsistent
with iPRA.
\end{corollary}
\begin{proof}
The term and formula languages, ordinary readings, and designated
contradiction have not changed. Apply
Theorem~\ref{rev:equivalence} to every sequent used in the existing
native-to-iPRA faithfulness and round-trip theorems. This transfers the
theorems \texttt{native\_iff\_pra}, \texttt{pra\_iff\_native},
\texttt{native\_roundTrip}, \texttt{pra\_native\_roundTrip}, and
\texttt{native\_inconsistent\_iff}; no new numerical representation
or quotient is introduced.
\end{proof}

\paragraph{Verification boundary.}
The argument above is a complete written derivation relative to the
supplied native kernel and its already established iPRA bridge. The
accompanying \texttt{Reversal.lean} is a proposed formalization using
the current source names and deeply embedded rules. It contains no
\texttt{sorry} or added axiom, but it has not been compiled in this
environment and is not claimed as machine verification. Checker-size
measurements are separate from this mathematical equivalence theorem.
